% Lecture notes template for the /lecture skill (fixed format).
% Replace every <...>. Keep only the sections this topic needs; copy macros from previous lectures.
\documentclass[11pt,a4paper]{article}
\usepackage[margin=0.8in]{geometry}
\usepackage{amsmath,amssymb,amsthm,bm}
\usepackage{booktabs}
\usepackage{enumitem}
\usepackage{xcolor}
\usepackage{titlesec}
\usepackage{fancyhdr}
% Only if the lecture has an algorithm:
% \usepackage{algorithm}
% \usepackage[noend]{algpseudocode}
\usepackage{hyperref}

% --- Colours ---
\definecolor{primary}{HTML}{D32F2F}
\definecolor{secondary}{HTML}{455A64}

\hypersetup{colorlinks, linkcolor=primary, urlcolor=primary,
    pdftitle={<Course code> - Lecture <N>: <Lecture title>}}

% --- Section titles ---
\titleformat{\section}{\large\bfseries\color{primary}}{\thesection}{1em}{}
\titleformat{\subsection}{\normalsize\bfseries\color{secondary}}{\thesubsection}{1em}{}

% --- Page style: running header, page number bottom right ---
\pagestyle{fancy}
\fancyhf{}
\lhead{\footnotesize Lecture <N>: <Short lecture title>}
\rhead{\footnotesize <Course code>: <Course title>}
\rfoot{\thepage}

\newtheorem{theorem}{Theorem}
\newtheorem{lemma}[theorem]{Lemma}
\theoremstyle{definition}
\newtheorem{definition}{Definition}
\newtheorem{example}{Example}
\theoremstyle{remark}
\newtheorem*{remark}{Remark}

\newcommand{\R}{\mathbb{R}}

\begin{document}

\begin{center}
    {\LARGE \bfseries \color{primary} <Course code>: <Course title>} \\[0.5em]
    {\large \bfseries Lecture <N>: <Lecture title>} \\[0.5em]
    {\small Instructor: <Instructor>}
\end{center}

\vspace{1em}

\section{<Motivation>}
% What students know, and the question this lecture answers. One or two paragraphs.

\section{<Core section>}
% Definitions, ideas, results with short proofs, worked examples, in strict order.
% Add as many core sections as the topic needs, and no more.

% \section{<The method>}   % only if the topic is a method: complete algorithm (+ cost if relevant)

% No summary section.
\section{Problems}   % unless COURSE.md turns problems off; typically 3-6, each with a purpose
\begin{enumerate}
\item
\end{enumerate}

\end{document}
